% =====================================================================
%  Do Television Series Decline Over Their Run?
%  Technical report -- compiles with pdfLaTeX on Overleaf as-is.
%  Upload the nine figures/fig0*.png files alongside this .tex.
% =====================================================================
\documentclass[11pt,a4paper]{article}

\usepackage[utf8]{inputenc}
\usepackage[T1]{fontenc}
\usepackage{lmodern}
% 3.4 cm side margins give a ~14 cm measure, about 72 characters a line. The
% earlier 2.4 cm setting produced a 16.2 cm measure and a median line of 92
% characters, well past the point where a reader starts losing their place.
\usepackage[a4paper,left=3.4cm,right=3.4cm,top=2.8cm,bottom=2.8cm]{geometry}
\usepackage{graphicx}
\usepackage{booktabs}
\usepackage{amsmath,amssymb}
\usepackage[font=small,labelfont=bf]{caption}
\usepackage{microtype}
\usepackage[colorlinks=true,linkcolor=black,citecolor=black,urlcolor=blue]{hyperref}
% Keeps each section's floats inside that section instead of letting them queue
% up and land three-to-a-page with no text between them.
\usepackage{placeins}
% \captionof outside a float (appendix table).
\usepackage{capt-of}

% Find the figures whether they sit next to the .tex or in a folder. The paper/
% plates come first: same data as the essay figures, but with no in-image title
% block (the caption does that job), a white surface, and labels sized to be
% legible once scaled into this column.
\graphicspath{{./}{paper/}{figures/paper/}{../figures/paper/}{figures/}{../figures/}}

\title{\bfseries Do Television Series Decline Over Their Run?\\[0.3em]
\large Evidence from 192{,}720 IMDb Episode Ratings}
\author{Tar\i k \c{S}ahar\\[0.2em]
\small Istanbul Technical University, Data Science and Analytics}
\date{Data snapshot: 22 July 2026}

\begin{document}
\maketitle

\begin{abstract}
\noindent
The belief that television series decline over their run is widely held and
rarely measured. Using the complete IMDb episode-rating corpus, we construct a
panel of 192,720 rated episodes from 3,234 series and estimate a per-show
quality trajectory as the slope of a vote-weighted least-squares fit of episode
rating on normalised episode order. We find the folk belief is false for the
typical series: only 16.9\% of shows decline by at least half a rating point
across their entire run, 57.4\% are essentially flat, and 25.7\% clearly rise.
A coupling-free test---regressing each show's second-half mean on its
separately measured first-half mean---yields a correlation of $+0.793$ and a
slope of $0.886$ against a noise-only expectation of $0.983$, implying a real
but negligible extra decline of 0.07 rating points for the strongest starters.
Season finales exceed the remainder of their own season 72.6\% of the time, and
series finales exceed their final season 71.6\% of the time, so the moment the
belief expects collapse is typically a local maximum. Among ended series the
final season is weaker than the rest of the run in 50.5\% of cases with a
median change of $-0.006$, and only 12.6\% lose half a point or more; this
conclusion is insensitive to the threshold chosen. Decline is concentrated
rather than general: show length ($-0.026$ per season, $t=-6.8$) and premiere
year ($+0.048$ per decade, $t=+4.8$) each survive mutual adjustment, whereas
genre is largely uninformative once age and length are controlled. The model
explains only 6.0\% of between-show variance, so these are reliable but small
effects. Finally, the belief describes the series it is drawn from better than
the population: the clear-decline share rises monotonically with a series' vote
count, reaching 43.3\% among the 60 most-watched series against a 16.9\% base
rate, and this gradient survives controls for precision, length and era. Because
vote counts respond to disappointment as well as recording it, we report this as
descriptive; it offers a measured account of how a belief this durable coexists
with evidence this consistent.

\medskip
\noindent\textbf{Keywords:} television ratings, quality trajectories,
regression to the mean, availability bias, IMDb
\end{abstract}

\section{Introduction}

That television series deteriorate over time is close to received wisdom. The
claim is specific enough to be testable---it asserts that a show's quality
declines monotonically or near-monotonically with time on air---yet it is
usually supported by enumeration of memorable cases rather than by measurement.

This report tests the claim directly. We ask six questions:

\begin{enumerate}\itemsep2pt
  \item What is the distribution of per-show quality trajectories?
  \item Is a show's later quality predictable from its earlier quality, and
        does high early quality decay?
  \item Do finales---season and series---underperform, as the belief implies?
  \item Is there a ``final-season curse''?
  \item If decline exists, which show characteristics predict it?
  \item Does the belief describe the series it is actually drawn from---the
        widely-watched ones---better than it describes the population?
\end{enumerate}

Our contribution is not a new method but a disciplined application of standard
ones to a question usually argued anecdotally, with explicit attention to the
artefacts that make naive versions of this analysis produce the opposite
answer. Section~\ref{sec:robust} documents five such artefacts, one of which
produced an apparently strong effect that survived neither reformulation nor
replication.

\FloatBarrier
\section{Data}

\FloatBarrier
\subsection{Source}

All ratings come from the IMDb public datasets~\cite{imdb}, downloaded on
\textbf{22 July 2026}. IMDb republishes these files daily; counts drift as new
episodes accumulate ratings and borderline series cross inclusion thresholds,
so the exact figures below are reproducible only against the committed
snapshot. Three tables are used: \texttt{title.episode} (mapping each episode
to its parent series, season and episode number), \texttt{title.ratings} (mean
rating and vote count per title) and \texttt{title.basics} (title type,
primary title, start year, end year, genres).

\FloatBarrier
\subsection{Construction}

Episodes with a missing season or episode number are dropped, since ordering
within a series is required. The episode table is inner-joined to ratings---so
only episodes carrying at least one vote survive---and then to the series rows
of \texttt{title.basics} restricted to \texttt{tvSeries} and
\texttt{tvMiniSeries}. Episodes are ordered by season then episode number, and
a within-show sequential index is assigned. Series with an empty end year are
flagged as ongoing rather than dropped. The result is one row per rated
episode.

\FloatBarrier
\subsection{Inclusion criteria}

Six filters define the analysis population (Table~\ref{tab:incl}).

\begin{table}[htbp]
\centering
\small
\caption{Inclusion criteria applied to the joined episode table.}
\label{tab:incl}
\begin{tabular}{p{3.6cm}p{4.0cm}p{5.2cm}}
\toprule
\textbf{Criterion} & \textbf{Threshold} & \textbf{Rationale} \\
\midrule
Title type & \texttt{tvSeries}, \texttt{tvMiniSeries} & excludes films and specials \\
Rated episodes & $\ge 13$ & a trajectory needs enough observations \\
Seasons & $\ge 2$ & decline across seasons is otherwise undefined \\
Mean votes per episode & $\ge 50$ & limits noise from thinly-rated series \\
Genre exclusions & not Game-Show, Reality-TV, Talk-Show, News & non-narrative formats have no season arc to decline \\
Documentary cap & kept only if $<100$ episodes & separates season-based series from magazine strands \\
\bottomrule
\end{tabular}
\end{table}

The documentary rule is platform-independent, since the data carry no channel
or platform field.

\FloatBarrier
\subsection{Sample}

Table~\ref{tab:sample} describes the resulting panel.

\begin{table}[htbp]
\centering
\small
\caption{The analysis sample after all six filters.}
\label{tab:sample}
\begin{tabular}{lp{9.0cm}}
\toprule
\textbf{Property} & \textbf{Value} \\
\midrule
Series & 3,234 \\
Rated episodes & 192,720 \\
Total votes represented & 160,926,754 \\
Premiere years & 1949--2026 \\
Ongoing at snapshot & 524 (16.2\%) \\
Episodes per series & median 36 (IQR 22--69, max 1,234) \\
Seasons per series & median 3 (IQR 2--5, max 70) \\
Mean rating per series & median 7.72 (IQR 7.39--8.03) \\
Votes per episode & median 190 \\
\bottomrule
\end{tabular}
\end{table}

The sample is right-skewed in length: 485 series exceed 100 episodes and 127
exceed 200. These were retained rather than trimmed; the trajectory metric
normalises episode order to $[0,1]$ so that run length cannot mechanically
inflate a slope, and length is analysed as an explanatory variable in its own
right (Section~\ref{sec:cross}).

\FloatBarrier
\section{Methods}

\FloatBarrier
\subsection{Per-show trajectory}
\label{sec:traj}

For each series, let $x_i \in [0,1]$ be the normalised position of episode $i$
in the run and $y_i$ its mean rating. We fit
\[
  y_i = a + b\,x_i + \varepsilon_i
\]
by weighted least squares with weights $w_i = \sqrt{v_i}$, where $v_i$ is the
episode's vote count.

\begin{figure}[htbp]
\centering
\includegraphics[width=\linewidth]{fig08_example_trajectories.png}
\caption{Three fitted trajectories, one from each outcome band defined below.
Points are episode means sized by vote count; the straight line is the estimate
$\hat{b}$, the grey curve the quadratic of Section~\ref{sec:shape}.}
\end{figure}

The estimate $\hat{b}$ is the show's \emph{trend}: the rating points gained or
lost between first and last episode. A show is classified as \emph{clearly
declining} if $\hat{b} \le -0.5$, \emph{clearly rising} if $\hat{b} \ge +0.5$,
and \emph{essentially flat} otherwise. The half-point cut is a perceptibility
convention, not a statistical one; what estimation error costs it is examined
under \emph{Estimation error in the trend} in
Section~\ref{sec:robust}. (The \emph{Threshold sensitivity} paragraph there
sweeps the final-season cut of Section~\ref{sec:delta} instead, which is a
different quantity.)

\paragraph{Choice of weights.}
Vote counts vary substantially within a series (median max/min ratio 3.5,
exceeding 10 for a tenth of series), so an unweighted fit lets a 51-vote
episode influence the trend as much as a 500,000-vote one. Weighting by $v_i$
directly is optimal under an i.i.d.-votes noise model, but IMDb vote counts
violate it: notorious episodes attract votes for reasons unrelated to
measurement precision. In \emph{Game of Thrones}, four of the six most-voted
episodes belong to the final season and all are rated far below the series
mean; linear weighting nearly doubles the fitted slope relative to an
unweighted fit ($-2.206$ versus $-1.024$), while $\sqrt{v_i}$ gives $-1.546$.
Log-weighting was also examined and is nearly indistinguishable from unweighted
in practice. We therefore use $\sqrt{v_i}$ throughout.

\begin{figure}[htbp]
\centering
\includegraphics[width=\linewidth]{fig09_weighting_choice.png}
\caption{The same 73 episodes under three weighting schemes. Raw vote weighting
lets the heavily-voted final-season cluster set the trend for the whole run;
the single most-voted episode is in fact a mid-run 9.9.}
\end{figure}

\FloatBarrier
\subsection{Final-season effect}
\label{sec:delta}

The slope in Section~\ref{sec:traj} answers ``does this series decline at a
roughly constant rate'', which is \emph{not} the final-season question: a
straight line fitted across all episodes is anchored by the majority and
understates a late, sharp drop. For \emph{Game of Thrones} the true season-8
mean is 6.40 against 8.95 for seasons 1--7, yet the fitted line predicts
7.6--8.3 in that range under every weighting scheme. We therefore define, for
ended series only,
\[
  \Delta = \bar{y}_{\text{final season}} - \bar{y}_{\text{all earlier episodes}},
\]
and use $\Delta$ for all final-season claims. A series is termed \emph{clearly
cursed} if $\Delta \le -0.5$.

\FloatBarrier
\subsection{Trajectory shape}
\label{sec:shape}

A weighted quadratic $y = a + bx + cx^2$ is fitted per series with the same
weights. Where the vertex $-b/2c$ lies inside $[0.15, 0.85]$ and $|c| > 0.05$,
the series is labelled \textbf{jumped-the-shark} ($c<0$: interior maximum) or
\textbf{found-itself} ($c>0$: interior minimum); otherwise \textbf{straight
line}, meaning the fitted curve has no turning point inside the observed run
and is monotone over it. Note this last class is not ``flat'': it contains
steep monotone decliners and risers alike. Volatility is the weighted RMSE of
the linear fit's residuals.

\FloatBarrier
\subsection{Peak location}

For each series we take the argmax of season means, reported both as an
absolute season number and as a normalised rank
$(\mathrm{rank}-1)/(n_{\text{seasons}}-1)$. Because the sample is dominated by
short series, aggregate statements about the peak are reported stratified by
season count.

\FloatBarrier
\subsection{Cross-sectional model}
\label{sec:model}

To separate era, length and genre we estimate by OLS
\[
  \hat{b}_j = \beta_0 + \beta_1 \mathrm{year}_j + \beta_2 \mathrm{seasons}_j
  + \beta_3 \mathrm{ongoing}_j + \sum_g \gamma_g \mathbf{1}[g \in \mathrm{genres}_j] + u_j
\]
over all $n = 3{,}234$ series, with premiere year centred and genres entered as
16 non-exclusive indicators. Classical standard errors are reported; given the
sample size and the magnitude of the leading $t$-statistics, robust
alternatives would not change any conclusion drawn here.

\FloatBarrier
\subsection{Software and figures}

The dataset is built and analysed with NumPy~\cite{numpy} and
pandas~\cite{pandas}; figures are drawn with Matplotlib~\cite{matplotlib}. Each
figure recomputes its own inputs from the processed table and prints the values
it annotates, so a plate cannot drift away from the number quoted beside it.
Figure colours were checked rather than chosen by eye: the palette was
validated for colour-vision deficiency using the Machado \emph{et
al.}~\cite{machado} simulation at full severity, which rejected one pair of
hues that had been used together.

\FloatBarrier
\section{Results}

\FloatBarrier
\subsection{The distribution of trajectories}
\label{sec:dist}

\begin{figure}[htbp]
\centering
\includegraphics[width=\linewidth]{fig01_slope_distribution.png}
\caption{Distribution of the per-show trend $\hat{b}$ across all $3{,}234$ series. Red marks clear decline ($\hat{b} \le -0.5$), blue clear rise, grey the flat band between them; the black rule is the median. The axis is clipped to $\pm 3$ points, which places 18 series (0.6\%) outside it.}
\end{figure}

The trend distribution is centred slightly above zero (mean $+0.073$, median
$+0.109$, SD $0.778$, IQR $-0.296$ to $+0.515$). By sign, 57.4\% of series
trend upward and 42.6\% downward. By magnitude, \textbf{16.9\% decline
clearly, 57.4\% are essentially flat, and 25.7\% rise clearly}. Under the folk
belief we would expect a distribution massed below zero; we observe a
distribution massed \emph{at} zero with a mild positive tilt.

\FloatBarrier
\subsection{Stability across a series' run}
\label{sec:halves}

\begin{figure}[htbp]
\centering
\includegraphics[width=\linewidth]{fig02_halves_stability.png}
\caption{Each series' first-half mean against its separately measured second-half mean, all $3{,}234$ series. The grey line is equality, the orange line the fitted relationship. Right: mean change from first half to second, by quartile of first-half rating.}
\end{figure}

Splitting each series into disjoint halves, first-half and second-half means
correlate $\mathbf{+0.793}$. Regressing second on first yields a slope of
$\mathbf{0.886}$, below unity and therefore consistent with some regression to
the mean. The relevant benchmark is not 1 but the reliability of a half-mean:
an odd/even split-half correlation with the Spearman--Brown
correction~\cite{spearman,brown} gives $\mathbf{0.983}$, which is what a pure
``stable quality plus measurement noise'' process would produce. The observed
0.886 falls modestly below this, so a genuine extra decline of strong starters
exists---but it is very small. In levels, top-quartile starters (first-half
mean 8.32) lose \textbf{0.07} rating points in their second half, and
bottom-quartile starters (6.94) gain 0.07.

\FloatBarrier
\subsection{Finale structure}

\begin{figure}[htbp]
\centering
\includegraphics[width=\linewidth]{fig03_finale_premium.png}
\caption{Finale and premiere premia relative to the rest of the season. Left:
mean premium in rating points. Right: share of cases beating the comparison
group, against the 50\% line. Seasons of at least four rated episodes,
$n = 12{,}931$ seasons and $3{,}082$ series finales.}
\label{fig:finale}
\end{figure}

Restricting to seasons of at least four rated episodes ($n = 12{,}931$ seasons),
the season finale beats the rest of its own season by $+0.253$ rating points on
average and does so 72.6\% of the time; the season premiere is slightly below
its own season ($-0.040$, above baseline 43.3\%); and the series finale beats
its final-season body by $+0.256$, 71.6\% of the time
($n = 3{,}082$ series). Figure~\ref{fig:finale} plots all three.

Finales are systematically the strongest episode of their season, premieres are
mildly below average, and the asymmetry is large. Only 9.4\% of series finales
fall 0.5 points or more below their final season---the memorable ``finale
flop'' is a one-in-ten event.

\FloatBarrier
\subsection{The final-season effect}

\begin{figure}[htbp]
\centering
\includegraphics[width=\linewidth]{fig04_final_season_curse.png}
\caption{Distribution of $\Delta$ for the $2{,}710$ ended series, with a sensitivity sweep over the ``cursed'' threshold. The black rule marks $\Delta = -0.5$, which 12.6\% of ended series fall below. The axis is clipped to $\pm 3$ points, leaving 6 series outside it.}
\end{figure}

Across 2,710 ended series, 50.5\% have a weaker final season and 49.2\% a
stronger one (ten end exactly level), with median $\Delta = -0.006$ and mean
$-0.049$. \textbf{12.6\%}
are clearly cursed at the $-0.5$ cut. The cursed minority is dominated by
well-known series (\emph{House of Cards} $-4.24$, \emph{Game of Thrones}
$-2.55$, \emph{Master of None} $-2.54$, \emph{The Promised Neverland}
$-2.28$), which suggests the mechanism sustaining the belief: the observable
instances are disproportionately the famous ones, so the cases available to
recall are not a fair sample of the cases that exist~\cite{tversky}.

That is measurable rather than merely illustrative. Cutting $\Delta$ by series
popularity gives a monotone gradient---11.6\% clearly cursed among series under
5,000 votes, then 12.5\%, 14.4\%, and \textbf{23.8\%} among the 42 household
names above 500,000, with median $\Delta$ moving from $+0.063$ to $-0.209$. The
top tier is small and its interval correspondingly wide (95\% Wilson
score~\cite{wilson} 13.5--38.5\%), so this is the weaker of the two fame
gradients we measure; the
slope-based one in Section~\ref{sec:pop} is sharper and better powered. Both
carry the endogeneity caveat stated there.

$\Delta$ and $\hat{b}$ correlate $+0.80$, so roughly a third of their variation
is independent and they are not interchangeable. Notably, the pure ``flat then
cliff'' case is essentially absent: all 32 series with $\Delta \le -1.0$ and at
least 30,000 votes also have a negative overall slope, which is
mechanical---a late crash necessarily tilts a fitted line downward.

\FloatBarrier
\subsection{Era, length and genre}
\label{sec:cross}

\begin{figure}[htbp]
\centering
\includegraphics[width=\linewidth]{fig05_era_length_genre.png}
\caption{OLS coefficients with 95\% confidence intervals, $n = 3{,}234$. A coloured dot clears its panel's significance bar, grey does not. The genre panel screens 16 coefficients simultaneously and so uses a Bonferroni-corrected bar ($|t| > 2.96$); the era/length panel tests pre-specified single hypotheses and uses $|t| > 1.96$. Under Benjamini--Hochberg, Documentary, Romance and Thriller would also be coloured; Biography would not. Era is shown per decade for comparability with the per-season length effect; the model itself is per year. Full table in Appendix~\ref{app:ols}.}
\end{figure}

Raw cross-sectional cuts are heavily confounded: median seasons per series
falls from 5 in the 1950s to 2 in the 2020s while the ongoing share rises from
3\% to 41\%. After mutual adjustment (full table in Appendix~\ref{app:ols}):

\begin{itemize}\itemsep2pt
  \item \textbf{Length:} $-0.0261$ per season ($t = -6.78$). Ended
        2--3-season series average $+0.217$; ended 6+-season series $-0.120$,
        with 60.8\% of 6+-season series declining.
  \item \textbf{Era:} $+0.0048$ per year, i.e. $+0.048$ per decade
        ($t = +4.81$). The effect survives restriction to ended series and to
        fixed length bands, so it is not a length or ongoing-status artefact.
        59.3\% of pre-2000 series decline against 38.6\% of later ones.
  \item \textbf{Genre:} six of the sixteen genre indicators reach significance
        individually, but sixteen simultaneous tests are expected to produce
        about 0.8 spurious flags at $\alpha = 0.05$, so the family requires a
        correction. Only \textbf{Animation} ($+0.233$, $t = 5.57$) and
        \textbf{Adventure} ($-0.165$, $t = -3.75$) survive Bonferroni
        ($\alpha = 0.0031$, $|t| > 2.96$). Benjamini--Hochberg~\cite{benjamini},
        which controls the false discovery rate and is the more appropriate
        screen for a set of exploratory indicators, additionally retains
        Documentary
        ($-0.235$), Romance ($-0.143$) and Thriller ($-0.153$).
        \textbf{Biography} ($-0.244$, $t = -2.25$) survives neither and is not
        reported as an effect: it is the weakest of the raw six and rests on
        53 series. Drama, Comedy, Crime and Action---the bulk of the
        sample---are indistinguishable from zero under any threshold.
\end{itemize}

\textbf{The model explains 6.0\% of between-show variance} ($R^2 = 0.0603$,
residual SD 0.756). The effects are estimated precisely but are small relative
to the idiosyncratic differences between series; no combination of a show's
age, length and genre comes close to determining its trajectory.

\FloatBarrier
\subsection{Shape and peak location}

\begin{figure}[htbp]
\centering
\includegraphics[width=\linewidth]{fig06_trajectory_shapes.png}
\caption{Share of series by trajectory shape, and the mean fitted arc of each class after centring on the series' own mean rating. The straight-line class contains both rising and falling series, so its average is their net. Class shares move by less than one point across curvature cut-offs from 0 to 0.20; roughly a fifth of the found-itself class rests on thinly-voted late episodes (Section~\ref{sec:robust}).}
\end{figure}

The rise-then-fall arc is the single most common shape (\textbf{44.6\%}),
against 35.2\% straight-line and 20.3\% found-itself. Within the
jumped-the-shark class the vertex sits at a median normalised position of 0.50,
with only 6.2\% of peaks in the 0.15--0.25 band, so the label describes a
genuine interior maximum rather than a decline from the outset. The arc is
nonetheless shallow: median rise to the peak is 0.37 rating points and median
fall from it 0.33. The classification is insensitive to the curvature
cut-off---class shares move from 44.6/35.1/20.3 at a threshold of 0 to
43.7/36.6/19.8 at 0.20---because the binding condition is the location of the
vertex, not the size of $c$.

\begin{figure}[htbp]
\centering
\includegraphics[width=\linewidth]{fig07_peak_location.png}
\caption{Location of the best season. Left: mean best season by run length, all $3{,}234$ series. Centre: share of series peaking in each third of their run, restricted to the 915 series with at least five seasons, where thirds are meaningful. Right: how often the first and the last season is the single best, all series.}
\end{figure}

The mean best season rises with run length (1.6 at two seasons, 2.5 at four,
3.1 at five, 3.8 at seven, 5.8 at eight or more) while its normalised position
remains near the middle (0.46--0.59 across all lengths). An earlier cut on
three-season series alone suggested a ``sophomore surge''; that is an
aggregation artefact of a sample dominated by short series, in which season 2
\emph{is} the middle. Importantly, the mean masks the distribution: among
series with at least five seasons ($n = 915$), the peak falls in the first
third 41.7\% of the time, the middle third 24.9\%, and the last third 33.3\%
---the middle is the \emph{least} likely location, and the mean sits at 0.5
only because the ends balance. Finally, the last season is the single best
season 37.2\% of the time against 27.1\% for the first, contradicting the
``first season is always best'' claim in aggregate.

\FloatBarrier
\subsection{Popularity and the provenance of the belief}
\label{sec:pop}

Section~\ref{sec:dist} is a statement about the average series. The belief under
test is not formed from the average series: it is formed from the few dozen
titles a mass audience has actually watched. Whether those behave like the
population is therefore a substantive question rather than a rhetorical one, and
it is measurable using each series' total vote count---the only popularity proxy
the data carries (Table~\ref{tab:pop}).

\begin{table}[htbp]
\centering
\small
\caption{Trend by series popularity. Intervals are 95\% Wilson score
intervals~\cite{wilson} for the clear-decline share; the top tier is small
enough that a bare percentage would overstate its precision.}
\label{tab:pop}
\begin{tabular}{lrrrr}
\toprule
\textbf{Tier} & \textbf{Series} & \textbf{Clear decline} & \textbf{95\% CI} &
\textbf{Median $\hat{b}$} \\
\midrule
Under 5,000 votes  & 1,184 & 14.8\% & 12.9--16.9 & $+0.224$ \\
5,000--50,000      & 1,545 & 15.1\% & 13.4--17.0 & $+0.078$ \\
50,000--500,000    &   445 & 24.7\% & 20.9--28.9 & $-0.041$ \\
Over 500,000       &    60 & \textbf{43.3\%} & 31.6--55.9 & $\mathbf{-0.392}$ \\
\bottomrule
\end{tabular}
\end{table}

The gradient is monotone, and the most-watched tier declines at roughly three
times the 16.9\% base rate---an interval that does not overlap the least-watched
tier's. Three attempts to dissolve it all fail.

\textbf{Not a precision artefact.} A noisier fit crosses a fixed $\pm 0.5$ cut
more often by chance, so systematically noisier estimates among obscure series
could manufacture the whole gradient. Median slope standard errors are
essentially flat across the four tiers (0.207, 0.174, 0.191, 0.205): popular
series have less noise per episode but genuinely bumpier ratings, and the two
roughly cancel. The intuition that popular means better-measured is wrong here,
which is why it was checked rather than assumed.

\textbf{Not the length effect in disguise.} Popular series run longer ($r =
0.37$ between log votes and season count) and long series were already shown to
decline. Holding the season band fixed, the gap survives in every band: 13.1\%
against 31.0\% at 2--3 seasons, 15.1\% against 36.6\% at 4--5, and 23.2\%
against 38.2\% at 6+ (obscure $<$ 20,000 votes against popular $\ge$ 200,000).

\textbf{Survives era and length jointly.} Regressing $\hat{b}$ on
$\log_{10}$(votes), season count and premiere year gives $-0.134$ per tenfold
increase in votes ($t = -6.2$), larger in magnitude than a whole extra season
($-0.019$) and estimated as sharply.

\textbf{Vote count is endogenous, and no causal claim is available.} This
project documented the mechanism itself while choosing a weighting scheme
(Section~\ref{sec:traj}): a final-season backlash inflates vote counts, which is
why four of \emph{Game of Thrones}' six most-voted episodes belong to season 8.
Decline raises votes as surely as votes track decline, so the arrow cannot be
pointed, and fame cannot be said to cause anything here.

What survives is descriptive. The series the
belief is drawn from do not behave like the typical series, and ``fewer than one
in six'' is not a statement about them. This is offered as a reconciliation
rather than a retraction: Section~\ref{sec:dist} remains the correct answer to
``do series decline'', and this section explains why an honest viewer with a
different sample would answer differently.

\FloatBarrier
\section{Robustness and threats to validity}
\label{sec:robust}

\paragraph{Mathematical coupling.}
An early formulation correlated each show's slope with its own fitted
intercept, producing $r = -0.373$ and the apparent finding that high-quality
series decline hardest. This is an artefact: both quantities are estimated from
the same episodes and share noise, so the correlation is partly
manufactured~\cite{barnett}. The disjoint-halves design in
Section~\ref{sec:halves} removes the shared-noise channel entirely and reduces
the effect to 0.07 rating points. We report the clean estimate and retain the
biased one only as a documented negative result.

\paragraph{Metric--question mismatch.}
As shown in Section~\ref{sec:delta}, a single linear slope cannot represent a
flat-then-cliff trajectory and understates late crashes. Using the slope as a
proxy for the final-season question would have understated \emph{Game of
Thrones}' season-8 drop by roughly 1.2--1.9 rating points. All final-season
claims therefore use the direct season contrast.

\paragraph{Threshold sensitivity.}
\label{sec:thresh}
The clearly-cursed share is 29.2\% at $\Delta \le -0.2$, 22.4\% at $-0.3$,
12.6\% at $-0.5$, 6.8\% at $-0.75$ and 3.7\% at $-1.0$: the conclusion that
cursed final seasons are a minority holds at every reasonable cut. As a
scale-free check we also computed $\Delta$ divided by the SD of the series' own
non-final-season episode ratings; 23\% of ended series reach a ``large'' effect
in Cohen's conventional sense~\cite{cohen}, and 98\% of series flagged by the
raw $-0.5$ rule also clear the large-effect bar, so the intuitive and
principled criteria select nearly the same set.

\paragraph{Low-vote tails.}
Within a series, episode position correlates $-0.74$ with log vote count: later
episodes are systematically less voted. Any late-leaning result is therefore on
weaker ground than an early-leaning one. This matters most for the found-itself
class: the median last-third episode carries only about 132 votes and 39.7\% of
found-itself series have a last third averaging under 100 votes. Re-fitting
after dropping sub-100-vote episodes, 56.5\% of found-itself series retain the
label, 20.2\% revert to decline or flat, and 23.4\% become untestable. The
label is not a weighting artefact (94.5\% survive an unweighted refit) but
roughly a fifth of it is a thin-data artefact, and we treat it accordingly.

\paragraph{Confounding.}
``Newer series rise'', ``shorter series rise'' and ``older series decline'' are
not three findings: newer series are shorter. The specification in
Section~\ref{sec:model} holds each fixed, and both era and length survive. We
additionally verified the era effect within ended series and within a fixed
2--3-season band, where it does not fade.

\paragraph{Estimation error in the trend.}
\label{sec:esterr}
Each $\hat{b}$ is an estimate, and Section~\ref{sec:dist} sorts estimates into
bands cut at $\pm 0.5$, so a series near a cut-off can land on the wrong side of
it by chance. Two questions follow, and they have opposite-looking answers.

\emph{Individual labels are soft.} The median slope standard error is 0.189 (IQR
0.139--0.255), not small against a band half-width of 0.5. Of the 545 series
labelled clearly declining, 285 (52\%) have a 95\% interval that still touches
$-0.5$; of the 832 labelled clearly rising, 510 (61\%) do. Only 55\% of series
have a slope distinguishable from zero at all. The bands are a reporting
convention over a continuous quantity, not a per-series verdict, and we make no
claims about individual series on their basis.

\emph{The population split is barely distorted.} Estimation errors are
independent of the true slopes, so the observed variance decomposes: $0.6049$
observed $= 0.5434$ real between-series variance $+\ 0.0614$ mean squared
estimation error, i.e. only 10\% of the spread between series is noise. This step
requires no assumption about the shape of either distribution and is the firmer
of the two results here.

\emph{Recovering a noise-free split does require a shape assumption, and the
obvious tools fail.} Two were tried and rejected, and we report them because they
explain why the estimate below is stated cautiously. Empirical-Bayes
shrinkage~\cite{efron} is the wrong instrument: posterior means are deliberately
over-shrunk as a set---here their variance is $0.4238$ against a true
$0.5434$, 22\% too narrow---so counting how many fall past a fixed cut
understates both tails by construction, reporting 15.2\% decline for a mechanical
rather than an empirical reason. A normal model for the true slopes is rejected by
the data: read parametrically, true $\sim N(0.073, 0.543)$ implies 21.9\%
decline, higher than the observed 16.9\% and thus the opposite sign, and it
predicts an \emph{observed} share of 23.1\% against the 16.9\% actually seen, so
it cannot reproduce the data it is fitted to. The slopes are skewed ($-0.64$) and
fat-tailed (excess kurtosis $+5.60$; 0.53\% lie beyond four SD against 0.006\%
under a normal).

What we use instead is a scale-family deconvolution assuming only normal
estimation \emph{error}---already assumed by every standard error in this
report---and no shape for the true slopes beyond preservation under scaling.
Solving for the scale $c$ at which the expected share of noisy estimates past the
cut, $\overline{\Phi((\text{cut}-t_i)/\hat\sigma_i)}$, matches the observed share
gives $c = 0.927$ and a corrected split of \textbf{15.6\% / 60.4\% / 23.9\%}. It
is fitted on the decline share alone and predicts a rise share of 25.0\% against
the 25.7\% observed, a genuine out-of-sample check.

\emph{The direction is robust; the magnitude is not, and only the direction is
claimed.} Noise smooths a density that is peaked inside the middle band, so mass
necessarily flows outward across that band's two edges---the middle is the only
class with two edges to lose across. This is visible without any deconvolution:
treating the estimates themselves as the truth and adding their own noise predicts
18.1\% observed decline against the 16.9\% seen. Estimation error therefore
inflates the apparent decline share by roughly a point, and we report the
uncorrected figures throughout, which is the conservative choice.

\FloatBarrier
\section{Limitations}

\paragraph{Ratings are not quality.}
IMDb scores are contributed by self-selected voters, not a representative
audience panel. Some low outliers reflect coordinated review-bombing rather
than dispersed judgement; vote weighting dampens but does not remove this.

\paragraph{Ratings are a present-day snapshot.}
Each rating is the value standing today, not the reception recorded at
broadcast. A series that ended in 2010 is scored in 2026 largely by viewers who
arrived later and already knew its ending. The finale premium may therefore
partly reflect retrospective softening. The dataset carries no time dimension
with which to test this.

\paragraph{Causes are unobserved.}
The design measures where ratings move, never why. Off-screen shocks are
indistinguishable from creative decline: \emph{House of Cards} lost its lead
actor amid scandal and had its final season cut from thirteen episodes to
eight; \emph{Top Gear} lost its presenting team; \emph{Scrubs} became a
different programme in its ninth season. All three appear in the data simply as
``decline''.

\paragraph{Format leakage.}
The genre filters remove the main non-narrative categories, but sketch and
review programmes tagged only as Comedy survive---most consequentially
\emph{Saturday Night Live} (1,010 episodes). A curated blacklist was rejected
as unreproducible; the affected count is too small to move any reported figure,
and we disclose it rather than patch it.

\paragraph{Survivorship in the vote threshold.}
Requiring at least 50 mean votes per episode excludes obscure series, so
results generalise to series with a measurable audience rather than to all
television.

\FloatBarrier
\section{Conclusion}

The proposition that television series decline over their run is not supported
for the typical series. Most series are close to flat across their entire run,
more rise than fall, later quality is strongly predictable from earlier
quality, and the extra decline of strong starters is an order of magnitude
smaller than the belief implies. Where the belief predicts collapse most
confidently---the finale---the data show a local maximum in roughly seven cases
out of ten.

Decline is nonetheless real in identifiable subsets: long-running series, older
series, and a 12.6\% minority of ended series whose final season drops
materially. It is also real in the subset the belief is actually drawn from.
Among the sixty most-watched series in the corpus, 43.3\% clearly decline
against a 16.9\% base rate, and the gradient across popularity tiers is monotone
and robust to controls for precision, length and era
(Section~\ref{sec:pop}). This is the reconciliation the paper ends on, and it is
not a retraction: the belief is a poor description of television and a fair
description of the television its holders have seen. The counterexamples are, by
construction, the series nobody talks about.

The effects we can identify are precisely estimated but jointly explain 6\% of
the variation between series. Whatever determines a show's trajectory is mostly
not its age, its length, or its genre.

\begin{thebibliography}{9}
\small

\bibitem{imdb}
IMDb. \emph{IMDb Non-Commercial Datasets.} \url{https://datasets.imdbws.com/}
Accessed 22 July 2026.

\bibitem{spearman}
Spearman, C. (1910). Correlation calculated from faulty data.
\emph{British Journal of Psychology}, 3(3), 271--295.

\bibitem{brown}
Brown, W. (1910). Some experimental results in the correlation of mental
abilities. \emph{British Journal of Psychology}, 3(3), 296--322.

\bibitem{barnett}
Barnett, A. G., van der Pols, J. C., \& Dobson, A. J. (2005). Regression to the
mean: what it is and how to deal with it. \emph{International Journal of
Epidemiology}, 34(1), 215--220.

\bibitem{cohen}
Cohen, J. (1988). \emph{Statistical Power Analysis for the Behavioral Sciences}
(2nd ed.). Lawrence Erlbaum Associates.

\bibitem{tversky}
Tversky, A., \& Kahneman, D. (1973). Availability: a heuristic for judging
frequency and probability. \emph{Cognitive Psychology}, 5(2), 207--232.

\bibitem{benjamini}
Benjamini, Y., \& Hochberg, Y. (1995). Controlling the false discovery rate: a
practical and powerful approach to multiple testing. \emph{Journal of the Royal
Statistical Society: Series B}, 57(1), 289--300.

\bibitem{wilson}
Wilson, E. B. (1927). Probable inference, the law of succession, and statistical
inference. \emph{Journal of the American Statistical Association}, 22(158),
209--212.

\bibitem{efron}
Efron, B., \& Morris, C. (1975). Data analysis using Stein's estimator and its
generalizations. \emph{Journal of the American Statistical Association},
70(350), 311--319.

\bibitem{numpy}
Harris, C. R., Millman, K. J., van der Walt, S. J., et al. (2020). Array
programming with NumPy. \emph{Nature}, 585(7825), 357--362.

\bibitem{pandas}
McKinney, W. (2010). Data structures for statistical computing in Python.
\emph{Proceedings of the 9th Python in Science Conference}, 56--61.

\bibitem{matplotlib}
Hunter, J. D. (2007). Matplotlib: a 2D graphics environment.
\emph{Computing in Science \& Engineering}, 9(3), 90--95.

\bibitem{machado}
Machado, G. M., Oliveira, M. M., \& Fernandes, L. A. F. (2009). A
physiologically-based model for simulation of color vision deficiency.
\emph{IEEE Transactions on Visualization and Computer Graphics}, 15(6),
1291--1298. Used to validate the figure palette for colour-vision deficiency.

\end{thebibliography}

\appendix

\section{Full cross-sectional model}
\label{app:ols}

OLS of the per-show trend on era, length, ongoing status and 16 non-exclusive
genre indicators. $n = 3{,}234$, $R^2 = 0.0603$, residual SD 0.756. Premiere
year is centred; the genre reference category is ``not carrying that genre''.
Table~\ref{tab:ols} reports every coefficient.

The sixteen genre coefficients are a family of simultaneous tests and the $t$
column must be read as such. Bonferroni ($\alpha = 0.05/16 = 0.0031$, i.e.
$|t| > 2.96$) leaves \textbf{Animation} and \textbf{Adventure};
Benjamini--Hochberg at the same level (reject $p \le 0.0129$) additionally leaves
Documentary, Romance and Thriller. Biography clears the raw 5\% bar but neither
correction, and should not be read as an effect. Era, length and ongoing status
are pre-specified single hypotheses, are not members of the genre family, and
keep the ordinary $|t| > 1.96$ threshold.

% A minipage rather than a bare center block: center CAN break across a page,
% and it did -- the first compile put this caption at the foot of one page and
% its table on the next. A \linewidth minipage cannot break, so the caption and
% the table always travel together.
\begin{minipage}{\linewidth}
\centering
\small
\captionof{table}{Full cross-sectional model. ``Series'' counts how many series
carry each genre tag; genres are non-exclusive.}
\label{tab:ols}
\begin{tabular}{lrrrr}
\toprule
\textbf{Variable} & \textbf{Coefficient} & \textbf{Std.\ error} & $t$ & \textbf{Series} \\
\midrule
Intercept                 & $0.2541$  & 0.0537 & $4.73$  & --- \\
Premiere year (per year)  & $0.0048$  & 0.0010 & $4.81$  & --- \\
Seasons                   & $-0.0261$ & 0.0038 & $-6.78$ & --- \\
Ongoing                   & $-0.0422$ & 0.0387 & $-1.09$ & 524 \\
\midrule
Animation                 & $0.2332$  & 0.0418 & $5.57$  & 680 \\
Action                    & $0.0414$  & 0.0413 & $1.00$  & 821 \\
Family                    & $0.0254$  & 0.0617 & $0.41$  & 172 \\
Comedy                    & $-0.0231$ & 0.0389 & $-0.60$ & 1,365 \\
Drama                     & $-0.0297$ & 0.0359 & $-0.83$ & 1,733 \\
Crime                     & $-0.0446$ & 0.0381 & $-1.17$ & 788 \\
History                   & $-0.0506$ & 0.0810 & $-0.62$ & 101 \\
Mystery                   & $-0.0552$ & 0.0467 & $-1.18$ & 405 \\
Sci-Fi                    & $-0.0726$ & 0.0745 & $-0.97$ & 114 \\
Horror                    & $-0.1064$ & 0.0733 & $-1.45$ & 125 \\
Fantasy                   & $-0.1086$ & 0.0566 & $-1.92$ & 214 \\
Romance                   & $-0.1428$ & 0.0531 & $-2.69$ & 254 \\
Thriller                  & $-0.1529$ & 0.0615 & $-2.49$ & 186 \\
Adventure                 & $-0.1651$ & 0.0440 & $-3.75$ & 716 \\
Documentary               & $-0.2345$ & 0.0858 & $-2.73$ & 101 \\
Biography                 & $-0.2436$ & 0.1083 & $-2.25$ & 53 \\
\bottomrule
\end{tabular}
\end{minipage}

\section{Reproduction}

Download the three IMDb tables into \texttt{data/raw/}, then run
\texttt{src/build.py}, the \texttt{phase2} scripts, and
\texttt{phase3\_figures.py} in that order. Full instructions are in the
repository README.

\end{document}
